\section{Conclusion}

A finite-radius PEC wire dipole was analytically initialized and numerically tuned for
operation near \SI{1}{GHz}. The first-order half-wave value of \SI{150}{mm} resonated
below the target in the documented CST model, while \SI{125}{mm} resonated above it.
The final \SI{135.6}{mm} configuration produced an exact CST marker at
\SI{0.99648}{GHz}, only \SI{0.352}{\percent} below the target, with
$S_{11}=-\SI{47.01186}{dB}$ relative to the selected \SI{73}{\ohm} reference.

The final far-field evidence shows the expected broadside toroidal pattern and axial
nulls. The displayed directivity of \SI{2.149}{dBi} and angular width of approximately
$78.4^\circ$ agree closely with the analytical thin half-wave values of
\SI{2.15}{dBi} and $78.08^\circ$. The verified conclusion is therefore that the preserved
CST model behaves as a tuned resonant wire dipole and reproduces canonical half-wave
radiation behavior.

The evidence remains analytical and simulated. A complete engineering qualification
would require the native model, raw exports, mesh and boundary convergence, a practical
feed or balun, fabrication, calibrated impedance measurement, and controlled far-field
testing.
